\section{Introduction}
Let $P\colon G= S_n \to \GL(n,\reals)$ be the 
standard permutation representation of 
the symmetric group~$S_n$ on $n$ letters.
The \emph{Birkhoff polytope} $B_n$ is by definition
the convex hull of all permutation matrices of size $n\times n$:
\[ B_n := \conv\{ P(\sigma) \mid \sigma\in S_n\}.
\]
In this note, we prove a conjecture of
Baumeister, Haase, Nill and Paffenholz~\cite[Conjecture~5.3]{BHNP09}
on the uniqueness of the Birkhoff polytope among
permutation polytopes.
In fact, we prove a slightly stronger result.

To state the result, we need the following notation.
Let $D\colon G\to \GL(d,\reals)$
be a representation over the reals.
The corresponding \emph{representation polytope},
$P(D)$, is the convex hull of the image of $D$:
\[ P(D):= \conv\{ D(g) \mid g\in G
               \}.
\]
If $D$ is a permutation representation, then the 
representation polytope is called a 
\emph{permutation polytope}.

Two representations 
$D_i\colon G_i \to \GL(d_i,\reals)$
(where $i=1$, $2$) are called 
\emph{effectively equivalent} if there is 
a group isomorphism $\phi\colon G_1\to G_2$
such that $D_1$ and $D_2\circ \phi$ are 
\emph{stably equivalent},
which means that $D_1$ and $D_2\circ \phi$ 
have the same nontrivial irreducible constituents 
(not necessarily occurring with the same multiplicities).
The representation polytopes of effectively representations are
affinely isomorphic
\cite[Theorem~2.4]{BaumeisterGrueninger15} \cite[\S~2]{BHNP09}.
The converse is not true, for example, 
when $D$ is the regular representation of a group,
then $P(D)$ is a simplex of dimension $\card{G}-1$.
Thus groups that are not even isomorphic as abstract groups,
may yield affinely equivalent representation polytopes.

From this viewpoint, the next result is somewhat surprising.
Recall that two polytopes 
$P$ and $Q$ are
\emph{combinatorially equivalent}
if there is a bijection between the vertices of $P$
and the vertices of $Q$ which maps faces of $P$
onto faces of $Q$.
Affinely equivalent polytopes are combinatorially equivalent,
but not conversely.

